
\chapter{QR-платежи и~кошельки}\label{ch:qr-wallets}
QR-код и~кошельковая обёртка не образуют отдельной платёжной системы --
за каждым из них стоит карточная сеть, A2A-контур или закрытый
оператор. Глава вводит общую таксономию методов оплаты, разбирает
три устойчивые модели QR-платежей, описывает российские кошельковые
сервисы и~фиксирует, почему карточные системы остаются центральным
методом, тогда как альтернативы выигрывают только в~нишах.

\section{Таксономия методов оплаты}\label{sec:payment-methods-taxonomy}

Карта, банковский перевод, кошелёк, рассрочка и цифровые деньги
устроены по-разному. Таксономия ниже сравнивает их по устройству,
а~не по маркетинговому описанию.

Первый класс -- карточные платежи, исторически самый
распространённый, в~который входят Visa, Mastercard, Мир, UnionPay
и AmEx. Списание инициирует продавец, отправляя
запрос по карточным реквизитам клиента и~получая либо авторизацию,
либо отказ. Сильные стороны -- глобальная инфраструктура, возвратный платёж
и~зрелая среда PCI DSS, слабая -- стоимость приёма.

Второй класс -- переводы со счёта на счёт (account-to-account, A2A),
где клиент сам подтверждает перевод со своего счёта на счёт продавца.
К A2A относятся СБП, SEPA Instant Credit Transfer, PIX и UPI~\cite{sbp-main,sepa-instant,pix,upi}.
Сильные стороны -- финальность и низкие комиссии, слабая -- более
узкая защита покупателя.

Третий класс -- кошельки, надстройка над базовой платёжной
инфраструктурой. Закрытый кошелёк (closed-loop wallet) хранит деньги
в~своей закрытой системе, как Alipay и WeChat Pay. Открытый кошелёк
(open-loop wallet) лишь оборачивает карточный инструмент токеном --
по этой модели работают Apple Pay и Google Pay. Для российских карт
после марта 2022~года оба сервиса недоступны, функциональный
эквивалент составляют Mir Pay и~СБПэй (см.~\ref{sec:russia-wallets}).

Четвёртый класс -- BNPL (Buy Now Pay Later), в~котором провайдер
оплачивает продавцу всю сумму сразу, а~клиент возвращает её частями по
собственному графику. Снаружи поток похож на эквайринг с~иной
комиссией, но экономическая модель другая, потому что в~центре стоит
кредитное решение провайдера, подробно разобранное в~главе~\ref{ch:bnpl}.

Пятый класс -- наличный сценарий (cash-based), в~котором платёж
начинается онлайн, но завершается офлайн через ваучер или платёжный референс.
Типичный пример -- OXXO Pay в~Мексике, где покупатель получает код
платежа и завершает операцию в~физической точке~\cite{oxxo-pay}.
Главный параметр здесь -- задержка подтверждения, потому что деньги
не приходят мгновенно.

Шестой класс -- цифровые деньги вне банковского депозита, внутри
которого выделяются две подгруппы. Цифровая валюта центрального
банка (central bank digital currency, CBDC) -- цифровая форма
обязательства центрального банка, и~российский цифровой рубль
разобран в~главе~\ref{ch:cbdc}. Криптовалюты и стейблкоины
строятся на иной модели эмиссии и регулирования, и~внешне их ставят
в~один ряд с~CBDC только потому, что ни те, ни другие не сводятся
к~обычному банковскому счёту.

\begin{table}[tbp]
  \begingroup
  \scriptsize
  \setlength{\tabcolsep}{2pt}
  \renewcommand{\arraystretch}{1.2}
  \makebox[\textwidth][l]{%
    \begin{minipage}{\fullwidth}
      \begin{tabularx}{\fullwidth}{
          @{}>{\RaggedRight\arraybackslash\bfseries}p{0.170\fullwidth}
          >{\RaggedRight\arraybackslash}p{0.075\fullwidth}
          >{\RaggedRight\arraybackslash}p{0.095\fullwidth}
          >{\RaggedRight\arraybackslash}p{0.090\fullwidth}
          >{\RaggedRight\arraybackslash}p{0.085\fullwidth}
          >{\RaggedRight\arraybackslash}p{0.145\fullwidth}
          >{\RaggedRight\arraybackslash}p{0.100\fullwidth}
          >{\RaggedRight\arraybackslash}p{0.155\fullwidth}@{}
        }
        \toprule
        \rowcolor{Panel}
        Метод оплаты & Тип & Финальность & Возвратный платёж & PCI
        scope & Международность & Скорость & Комиссия ТСП \\
        \midrule
        \rowcolor{Soft}
        Карты &
        \cellcolor{Bad!10}pull &
        \cellcolor{Bad!10}нет &
        \cellcolor{Bad!10}есть &
        \cellcolor{Bad!10}да &
        \cellcolor{Good!14}глобальный &
        \cellcolor{Good!14}мс &
        \cellcolor{Bad!10}высокая \\
        A2A (СБП/PIX) &
        \cellcolor{Good!14}push &
        \cellcolor{Good!14}да &
        \cellcolor{Good!14}нет &
        \cellcolor{Good!14}нет &
        \cellcolor{Ink!6}региональный &
        \cellcolor{Good!14}сек &
        \cellcolor{Good!14}низкая \\
        \rowcolor{Soft}
        Closed-loop кошелёк &
        \cellcolor{Ink!6}push/pull &
        \cellcolor{Ink!6}зависит &
        \cellcolor{Ink!6}ограничен &
        \cellcolor{Good!14}нет &
        \cellcolor{Bad!10}нишевый &
        \cellcolor{Good!14}мс &
        \cellcolor{Ink!6}средняя \\
        BNPL &
        \cellcolor{Bad!10}pull &
        \cellcolor{Bad!10}нет &
        \cellcolor{Bad!10}есть &
        \cellcolor{Ink!6}нет &
        \cellcolor{Ink!6}региональный &
        \cellcolor{Good!14}сек &
        \cellcolor{Ink!6}средняя \\
        \rowcolor{Soft}
        Cash-based &
        \cellcolor{Good!14}push &
        \cellcolor{Good!14}да &
        \cellcolor{Good!14}нет &
        \cellcolor{Good!14}нет &
        \cellcolor{Bad!10}нишевый &
        \cellcolor{Bad!10}часы / дни &
        \cellcolor{Ink!6}нет данных \\
        CBDC (цифровой рубль) &
        \cellcolor{Good!14}push &
        \cellcolor{Good!14}да &
        \cellcolor{Good!14}нет &
        \cellcolor{Good!14}нет &
        \cellcolor{Bad!10}региональный &
        \cellcolor{Ink!6}сек / мин &
        \cellcolor{Ink!6}нет данных \\
        \rowcolor{Soft}
        Crypto &
        \cellcolor{Good!14}push &
        \cellcolor{Good!14}да &
        \cellcolor{Good!14}нет &
        \cellcolor{Good!14}нет &
        \cellcolor{Ink!6}глобальный / нишевый &
        \cellcolor{Ink!6}мин / часы &
        \cellcolor{Ink!6}нет данных \\
        \bottomrule
      \end{tabularx}
      \par\smallskip
      \fcolorbox{Ink!18}{Soft}{%
        \parbox{\dimexpr\fullwidth-2\fboxsep-2\fboxrule\relax}{%
          \textcolor{Muted}{Зелёный цвет показывает сравнительное
            преимущество, красный указывает на ограничение, серый на нейтральную или
            зависящую от провайдера зону.
            * Данные оценочные и меняются в~зависимости от конкретного
          провайдера, маршрута и географии.}%
        }%
      }
    \end{minipage}%
  }
  \endgroup
  \caption{\enquote{Методы оплаты / параметры}}\label{fig:payment-methods-matrix}
\end{table}

\section{QR-платежи: три модели}\label{sec:qr-models}

Фраза \enquote{мы принимаем QR} без указания инфраструктуры не сообщает,
кто подтверждает платёж, где наступает финальность и кто разбирает
спорные операции. Сам по себе QR-код -- лишь способ передать платёжные данные без
NFC-оборудования или ручного ввода реквизитов, а~за ним стоит
A2A-перевод, закрытый кошелёк или карточная транзакция.

Физически QR-код в~платежах используется в~двух режимах,
которые отличаются тем, кто инициирует сканирование.
В режиме merchant-presented mode (MPM) продавец показывает код, а~клиент
сканирует его камерой телефона. В режиме customer-presented mode (CPM)
клиент открывает динамический код в~приложении, а~продавец считывает
его сканером на кассе. MPM дешевле в~развёртывании, потому что достаточно
напечатанной наклейки или экрана на кассе, тогда как CPM быстрее на высоком
потоке покупателей, но требует от продавца считывающее
оборудование. Обе модели существуют внутри каждой из трёх систем
ниже, однако пропорция их использования различается.

\subsection*{СБП: QR поверх A2A-перевода}

В~СБП QR-код ведёт в~мгновенный перевод со счёта клиента на счёт
ТСП через инфраструктуру НСПК~\cite{sbp-main}. Система работает
как push-платёж, потому что клиент инициирует и подтверждает списание сам.

Статический QR -- напечатанный код, который содержит
идентификатор ТСП, но не содержит сумму. Клиент сканирует код
в~приложении своего банка, вводит сумму вручную и подтверждает перевод.
Формат годится для малого бизнеса, поскольку не нужно кассовое ПО,
достаточно наклейки. Динамический QR генерируется кассовой системой
под конкретную покупку и~содержит реквизиты получателя, сумму
и назначение платежа. Клиент сканирует код, видит предзаполненные
данные и подтверждает одним действием, а~кассовая система получает
обратный вызов об успешном переводе за несколько
секунд~\cite{sbp-faq-business}.

Помимо QR, СБП поддерживает оплату по платёжной ссылке, через NFC
и~через привязку счёта в~приложении
СБПэй~\cite{sbpay-app,sbp-main}. Все эти каналы ведут к~одному
и~тому же A2A-переводу внутри инфраструктуры НСПК; различается
только способ, которым клиент передаёт согласие на списание.

Тарифы С2В-переводов в~СБП устанавливает Совет директоров Банка
России. Ставка межбанковского вознаграждения зависит от MCC ТСП
и укладывается в~диапазон от 0\% до 0{,}7\% от суммы
перевода. Для платежей в~пользу государства она равна 0\%, для услуг ЖКХ
ограничена 0{,}2\%, для социально значимых категорий
(продукты, лекарства, медицина и~т.\,п.) -- 0{,}4\%, для прочих
категорий -- 0{,}7\%~\cite{cbr-sbp-c2b-tariffs}. Для сравнения,
\textit{interchange} в~карточном эквайринге начинается от 1\% и превышает 2\%.
Регулируемая ставка -- осознанное решение Банка России снизить
стоимость приёма платежей для бизнеса.

Механизма возвратного платежа в~СБП нет, поскольку перевод финален
с~момента зачисления на счёт получателя. Если клиент хочет вернуть
деньги за покупку, продавец инициирует обратный перевод через свой
банк~\cite{sbp-faq-business}. Ошибочный перевод не тому получателю
также требует согласия получателя на возврат, если иной
порядок не установлен законом или решением
суда~\cite{cbr-sbp-faq-mistaken-transfer}. Эта модель выгоднее
продавцу, но смещает риск спора на покупателя -- прямая
противоположность карточному возвратному платежу.

\subsection*{Alipay и WeChat Pay: QR поверх закрытого
кошелька}\label{sec:qr-wallet-closed-loop}

В Alipay и WeChat Pay QR-код ведёт не в~банковский перевод,
а~в~платёж внутри закрытой системы кошелька. Средства пользователя
хранятся на балансе кошелька, пополняемом с~банковского счёта или
карты, и~списание происходит на стороне оператора кошелька без
выхода в~межбанковский клиринг.

Оба режима, MPM и CPM, здесь активно используются, причём
в~документации Alipay+ merchant-presented mode описан как основной
сценарий для сторонних ТСП. Продавец показывает QR-код,
клиент сканирует его в~приложении кошелька, оператор списывает
средства с~баланса клиента и зачисляет их на счёт
ТСП~\cite{alipayplus-mpm}. Customer-presented mode распространён
в~высокопоточных офлайн-точках вроде сетевых магазинов, автоматов
и общественного транспорта. Клиент показывает динамический штрихкод
из приложения, кассир сканирует его, и~оплата проходит без
дополнительного подтверждения со стороны клиента при сумме ниже
установленного оператором порога.

Оператор кошелька совмещает роли, которые в~карточной сети
разделены между эмитентом и эквайером. Он ведёт балансы и~клиентов,
и~ТСП, проводит расчёт внутри своего реестра и~сам определяет
правила диспутов. Для клиента весь платёжный опыт от пополнения
до возврата замкнут внутри одного приложения, а~для продавца
интеграция идёт через API оператора кошелька, а~не через банк-эквайер.

Трансграничное расширение работает через партнёрскую сеть, в~которой
Alipay+ объединяет локальные кошельки -- GCash, TrueMoney, KakaoPay
и~другие, -- и~их пользователи платят в~точках, принимающих Alipay, за
пределами домашней страны. Продавец подключается к~одному контракту
с~локальным партнёром Alipay+ и~получает доступ к~пользователям
нескольких кошельков~\cite{alipayplus-mpm}. Модель отличается
от карточной глобальности Visa или Mastercard, потому что вместо единого
глобального протокола работает федерация кошельков, каждый из которых
действует в~собственной расчётной инфраструктуре.

\subsection*{PIX: QR, встроенный в~национальную платёжную
инфраструктуру}\label{sec:qr-pix}

PIX -- система мгновенных платежей Banco Central do Brasil,
запущенная в~ноябре 2020 года~\cite{pix}. Как и~СБП, PIX строится
вокруг мгновенного A2A-перевода с~финальностью в~секундах и~низкой
комиссией. Расчёт происходит в Sistema de Pagamentos Instantâneos
(SPI) -- выделенной инфраструктуре мгновенных платежей,
управляемой центральным банком~\cite{pix-spi}.

Участие обязательно для всех финансовых институтов с~более чем 500
тыс.\ активных счетов. PIX отличается этим от европейских A2A-систем,
где подключение к~мгновенным платежам долгое время оставалось
добровольным, и~обязательность дала PIX покрытие с~момента запуска.

QR-код в PIX -- штатный интерфейс самой системы, а~не надстройка
PSP или банка, и~Banco Central do Brasil определяет два типа QR.
Статический многоразовый код содержит идентификатор получателя
и опционально сумму, а~динамический одноразовый генерируется под
конкретную транзакцию с~полным набором данных. Помимо сканирования
QR, пользователь применяет режим Pix Copia
e Cola -- текстовое представление тех же данных, которое копируется
и вставляется в~банковское приложение~\cite{pix-faq}. Режим
решает сценарий, в~котором и~платёж, и QR-код находятся на одном
устройстве и физическое сканирование камерой невозможно.

Комиссия для физических лиц за переводы и~платежи через PIX
равна нулю -- это установлено регуляторно. Для юридических лиц
тарифы определяются банком-участником, но значительно ниже
карточного эквайринга. Как и~в~СБП, возвратный платёж отсутствует,
поскольку PIX работает как push-платёж и перевод финален с~момента
зачисления на счёт получателя.

\subsection*{Сравнение трёх моделей}

\begin{table}[tbp]
  \begingroup
  \scriptsize
  \setlength{\tabcolsep}{3pt}
  \renewcommand{\arraystretch}{1.2}
  \makebox[\textwidth][l]{%
    \begin{minipage}{\fullwidth}
      \begin{tabularx}{\fullwidth}{
          @{}>{\RaggedRight\arraybackslash\bfseries}p{0.16\fullwidth}
          >{\RaggedRight\arraybackslash}p{0.27\fullwidth}
          >{\RaggedRight\arraybackslash}p{0.27\fullwidth}
          >{\RaggedRight\arraybackslash}p{0.27\fullwidth}@{}
        }
        \toprule
        \rowcolor{Panel}
        & СБП & Alipay / WeChat Pay & PIX \\
        \midrule
        \rowcolor{Soft}
        Инфраструктура & A2A-перевод (межбанковский) & Закрытый кошелёк & A2A-перевод (межбанковский) \\
        Оператор & НСПК (Банк России) & Ant Group / Tencent & Banco Central do Brasil \\
        \rowcolor{Soft}
        Расчёт & Через НСПК, секунды & Внутри реестра оператора & SPI при BCB, секунды \\
        Финальность & Мгновенная (push) & Мгновенная (внутри системы) & Мгновенная (push) \\
        \rowcolor{Soft}
        Возвратный платёж & Нет & Правила оператора & Нет \\
        Статический QR & Да & Да & Да \\
        \rowcolor{Soft}
        Динамический QR & Да & Да & Да \\
        Основной QR-режим & MPM & MPM + CPM & MPM \\
        \rowcolor{Soft}
        Комиссия С2B & 0\%--0{,}7\% (регулятор) & Устанавливает оператор & 0\% (физ.\ лица) \\
        Трансграничность & Региональный & Федерация кошельков & Региональный \\
        \bottomrule
      \end{tabularx}
    \end{minipage}%
  }
  \endgroup
  \caption{Три модели QR-платежей}\label{tab:qr-three-models}
\end{table}

Три модели используют один и~тот же физический канал -- QR-код, --
но за ним стоят разные системы с~разными свойствами (табл.~\ref{tab:qr-three-models}).
СБП и PIX проводят платёж через межбанковскую инфраструктуру, управляемую
центральным банком, тогда как кошелёк Alipay или WeChat Pay замыкает платёж
внутри собственного реестра. Первая архитектура даёт мгновенный
push-перевод с~регуляторно-ограниченной комиссией и отсутствием
возвратного платежа, вторая -- контроль оператора над всем циклом платежа,
включая правила диспутов, ценообразование и трансграничное расширение.

QR -- лишь упаковка, а~стоящая за ним платёжная инфраструктура определяет
финальность, стоимость, модель диспутов и географию покрытия.

\section{Российские платёжные приложения и кошельки}\label{sec:russia-wallets}

После марта 2022~года в~России сложилась собственная
структура платёжных приложений и кошельков. Она заменяет ряд
функций ушедших Apple Pay и Google Pay и одновременно отражает
специфику внутрироссийской инфраструктуры с~картами \enquote{Мир},
СБП и цифровым рублём. Ниже -- перечень основных решений, актуальных
на 2024--2026~годы. Конкретный набор поддерживаемых функций
и~состав банков-партнёров у~каждого сервиса меняется и сверяется
по официальной документации перед публикацией.

\textbf{Mir Pay} -- NFC-кошелёк на Android поверх токенизации НСПК
(см.~\ref{sec:mir-pay}). \textbf{СБПэй} -- мобильное приложение НСПК
поверх СБП с~оплатой по QR, NFC-табличке, кнопке и~ссылке за счёт
привязанного счёта, доступен и~на Android, и~на iOS
(см.~\ref{sec:sbp-sbpay}).

\textbf{SberPay} -- кошелёк и платёжный сервис ПАО \enquote{Сбербанк},
который поддерживает оплату на сайтах партнёров через кнопку SberPay
и~в~офлайн-точках с~поддержкой соответствующего сценария.
Для держателей карт \enquote{Мир} Сбербанка также реализованы
сценарии оплаты по QR-коду и NFC через мобильное приложение
банка~\cite{sber-pay}.

\textbf{T-Pay}, бывший \enquote{Tinkoff Pay}, -- платёжный
сервис АО \enquote{ТБанк}, бывшего Тинькофф Банка. Сервис работает
на сайтах партнёров и в~офлайн-точках, в~том числе с~использованием
NFC на Android, а~сценарии оплаты опираются на карты
банка и~счета клиентов в~его экосистеме~\cite{tpay}.

\textbf{Yandex Pay} -- платёжный сервис ООО \enquote{Яндекс.Пэй},
встраиваемый в~чекаут партнёрских сайтов и ориентированный на
безфрикционную оплату ранее сохранёнными картами. С 2024~года
он интегрирован с~сервисом рассрочки \enquote{Сплит}
(глава~\ref{ch:bnpl})~\cite{yandex-pay,yandex-split}.

\textbf{ЮKassa} -- платёжный сервис ООО НКО \enquote{ЮМани},
дочерней структуры Сбербанка. Работает прежде всего как PSP
для бизнеса с~приёмом карт, СБП и BNPL-сценариев. Часть
инфраструктуры делит с~одноимённым клиентским кошельком
\textbf{ЮMoney}, который сохраняет роль электронного средства
платежа для розничных пользователей в~продолжение традиции
Яндекс.Денег~\cite{yookassa,yoomoney}.

\textbf{Биоэквайринг и оплата по лицу} -- сценарий, который
НСПК и~Банк России развивают на основе Единой биометрической
системы (ЕБС). Покупатель оплачивает покупку, подтверждая
себя биометрией на специальном терминале, без карты и
телефона~\cite{nspk-biopay-overview,ebs-overview}.\footnote{Состав
торговых сетей, поддерживающих биоэквайринг, и операционные
сценарии вроде порога суммы без подтверждения и способов пополнения
ЕБС-привязки сверяются по сайту НСПК и~Банка России.}
В 2024--2026~годах сервис расширялся в~крупных розничных
сетях, и~хотя объёмы по сравнению с~СБП и~картами остаются
небольшими, направление становится самостоятельным каналом
оплаты.

\practicenote{%
  Для российского ТСП в 2024--2026 годах базовый чекаут
  обычно включает приём карт \enquote{Мир} (через эквайринг),
  СБП (по QR или кнопке), один или несколько кошельковых
  сценариев (SberPay, T-Pay, Yandex Pay), отдельный
  BNPL-сценарий для подходящих категорий и -- на отдельных
  типах ТСП -- биоэквайринг. Apple Pay и Google Pay
  для российских карт в~этот набор больше не входят,
  и~попытки их сохранить через зарубежную карту в~кошельке
  упираются в~антифрод и~в~ограничения схем.

}

\section{Структура способов оплаты в~России и~СНГ}\label{sec:payments-mix-cis}

Одна таблица с~процентами создаёт иллюзию единой структуры способов
оплаты по всему региону, тогда как устойчивыми остаются паттерны, а~не конкретные цифры.

Россия сочетает массовую карточную инфраструктуру и быстро растущие
альтернативы. Банк России указывает, что карты остаются самым
популярным средством безналичной оплаты, а~платежи по QR-коду
в 2024 году достигли 4{,}1~трлн рублей~\cite{cbr-results-2024-payments}.
Казахстан показывает другой профиль, потому что национальная статистика
отдельно выделяет крупный поток безналичных QR-платежей наряду
с~карточными операциями~\cite{nbk-payment-cards-2025jan}.
Узбекистан опирается на две национальные карточные системы --
HUMO и UZCARD~\cite{humo-faq,uzcard-about}. Армения сохраняет
собственную карточную инфраструктуру ArCa, и~оператор
системы указывает, что банки-участники выпускают карты ArCa
наряду с Visa, Mastercard и другими международными
системами~\cite{arca-about}.

СНГ -- не единый платёжный рынок, поэтому при локализации продукта нельзя
механически переносить российскую структуру способов оплаты на
Казахстан или армянский сценарий на Узбекистан без проверки локальной
инфраструктуры, роли QR и регуляторных требований.

\practicenote{%
  При локализации в~новой стране СНГ проверьте BIN-распределение
  трафика, наличие массового локального A2A-метода и требования к
  обязательному приёму локальных карт. Эти факторы обычно сильнее
  всего меняют экономику и конверсию.

}

\section{Почему карточные системы пока доминируют}\label{sec:card-dominance}

Карты сохраняют центральное место в~потребительских платежах из-за
наложения нескольких инерций.

Первая причина -- сетевые эффекты, потому что глобальные карточные
схемы уже выстроили массовую эмиссию, приём и привычный пользовательский
сценарий. Продавец, который принимает карты, автоматически доступен
большинству покупателей, а~новый метод должен одновременно получить
поддержку со стороны эквайера и~эмитента, и~это самый трудный вариант
сетевого эффекта.

Вторая причина -- международность, ведь карта, выпущенная в~одной стране,
работает во множестве других по единому протоколу, тогда как локальные
A2A-системы вроде СБП, PIX и UPI не взаимозаменяемы.

Третья причина -- защита покупателя через возвратный платёж, механизм
которого A2A-системы не воспроизводят без изменения самой идеи финальности.
Покупатель, оплативший картой незнакомому продавцу, имеет формальный
механизм возврата, а~покупатель, оплативший через A2A, чаще остаётся
с договорной защитой без схемного возврата.

Четвёртая причина -- инфраструктурная инерция, потому что мировая сеть
POS-терминалов уже выстроена под карточные схемы. Переоснащение
или добавление нового метода окупается только при массовом спросе, а
массовый спрос не возникает, пока продавцы не видят достаточного
пользовательского потока, и~это классический замкнутый круг внедрения.

Альтернативные методы выигрывают в~нишах -- низкая комиссия для
A2A, доступность для небанковских клиентов в~наличных сценариях,
условные платежи в~системе CBDC, рассрочка в BNPL и программируемые
трансграничные сценарии для криптоактивов. Глобальную
универсальность и инфраструктурный охват карт они пока не
воспроизводят, поэтому новые методы входят в~рынок рядом с
картами там, где их частное преимущество
перевешивает общее карточное удобство.

Если продавец работает только внутри России и~средний
чек ниже 100\,000~руб., СБП даёт минимальную комиссию
от 0 до 0{,}7\,\% в~зависимости от категории ТСП и мгновенное зачисление,
тогда как карты нужны как \textit{fallback} для клиентов без мобильного банка.
Если продавец работает международно, карты остаются основным
методом, а~A2A-системы других стран -- PIX или UPI India -- требуют
отдельной интеграции на каждый рынок. Если средний чек составляет
10\,000--150\,000~руб. и конверсия критична, BNPL повышает
конверсию за счёт рассрочки, но добавляет кредитный риск
на стороне BNPL-провайдера и~задержку в~расчёте. Если
покупатель не имеет банковского счёта, как в~части стран Азии
и Африки, единственный путь -- наличные через агентскую сеть или мобильные
деньги. Универсального метода не существует, и~оптимальный набор
определяется географией, средним чеком, профилем покупателя
и требованием к~механизму спора.
